\subsection{Lộ trình phát triển sản phẩm trong các giai đoạn tiếp theo}
Sản phẩm không chỉ là công cụ tạo host ảo, mà được định vị là một \textbf{AI Sales Operation Workflow} cho shop livestream nhỏ và SME: quản lý phiên live đa kênh, xử lý bình luận theo ngữ cảnh sản phẩm, hỗ trợ chốt đơn và tạo nội dung livestream bằng AI. Ở giai đoạn đầu, nhóm ưu tiên chứng minh năng lực lõi gồm AI Sales Agent, Knowledge Database, AI Video/Voice Generator và dashboard vận hành đa kênh. Khi sản phẩm đi vào thử nghiệm, trọng tâm chuyển sang đo độ chính xác của phản hồi, độ ổn định khi livestream, khả năng giảm tải nhân sự và mức sẵn sàng trả phí của nhóm khách hàng mục tiêu. Sau MVP, lộ trình mở rộng tập trung vào thương mại hóa tại TP.HCM, sau đó mở rộng ra các thành phố lớn và từng bước tích hợp sâu hơn với hệ thống quản lý bán hàng của khách hàng tổ chức.

\begin{center}
\scriptsize
\setlength{\tabcolsep}{3pt}
\renewcommand{\arraystretch}{1.28}
\begin{longtable}{|>{\raggedright\arraybackslash}p{2.0cm}|
    >{\raggedright\arraybackslash}p{3.0cm}|
    >{\raggedright\arraybackslash}p{4.0cm}|
    >{\raggedright\arraybackslash}p{3.35cm}|
    >{\raggedright\arraybackslash}p{2.65cm}|}
    \caption{Roadmap phát triển 1Stream từ R\&D đến mở rộng toàn quốc}
    \label{tab:roadmap-1stream} \\
    \hline
    \rowcolor{blue!15}
    \textbf{Giai đoạn} &
    \textbf{Trọng tâm chiến lược} &
    \textbf{Sản phẩm \& công nghệ} &
    \textbf{Thị trường \& vận hành} &
    \textbf{Chỉ số kiểm chứng} \\ \hline
    \endfirsthead

    \hline
    \rowcolor{blue!15}
    \textbf{Giai đoạn} &
    \textbf{Trọng tâm chiến lược} &
    \textbf{Sản phẩm \& công nghệ} &
    \textbf{Thị trường \& vận hành} &
    \textbf{Chỉ số kiểm chứng} \\ \hline
    \endhead

    \textbf{Giai đoạn 1}\newline \textbf{06--08/2026}\newline \textbf{Research \& Development} &
    Hoàn thiện nghiên cứu công nghệ lõi cho livestream bán hàng: LLM, xử lý
    bình luận realtime, TTS, AI video và quản lý phiên live đa kênh. &
    Xây dựng prototype gồm dashboard live preview, kho nội dung, auto-reply demo,
    Knowledge Database JSON/TF-IDF, phân loại intent, job video/TTS, micro-scene
    pipeline, FFmpeg và luồng livestream đa nền tảng qua RTMP. &
    Nghiên cứu hành vi livestream của shop nhỏ, seller có ít nhân sự và nhóm
    FMCG/beauty tại TP.HCM. Bootstrap bằng nguồn lực nội bộ; GPU thì có thể dùng nội bộ máy cá nhân ở giai đoạn đầu, API gọi free để thử nghiệm sản phẩm . & Có prototype chạy được cho live preview, RAG reply, sales-agent flow và
    video/TTS job; hoàn thành bộ test lõi cho intent, retrieval, order tools và
    micro-scene pipeline. \\ \hline

    \textbf{Giai đoạn 2}\newline \textbf{09--12/2026}\newline \textbf{Testing \& Validation} &
    Chạy thử nghiệm đóng để kiểm chứng nhu cầu, độ chính xác phản hồi và mức giảm
    tải vận hành cho người bán. &
    Hoàn thiện MVP v1: gom bình luận, phân loại ý định, trả lời theo dữ liệu sản
    phẩm, gợi ý/chốt đơn cơ bản và quản lý phiên live. Tích hợp AI Video/Voice
    vào MVP v2; kiểm thử WebSocket, job status và độ ổn định khi tạo video. &
    Alpha test với 5 shop/SME ngành FMCG
    hoặc beauty. Mở vòng pre-seed để tài trợ cloud, GPU, dữ liệu đánh giá và hỗ
    trợ khách hàng beta. &
    Giảm thời gian trả lời bình luận, giảm bỏ sót câu hỏi giá/ship/tồn kho, tỷ lệ
    câu trả lời cần người sửa ở mức chấp nhận được; có phản hồi rõ về setup và
    mức sẵn sàng trả phí. \\ \hline

    \textbf{Giai đoạn 3}\newline \textbf{01--04/2027}\newline \textbf{MVP Launch \& Market Entry} &
    Ra mắt MVP thương mại tại TP.HCM, ưu tiên micro-merchant và shop FMCG/beauty
    có tần suất livestream cao. &
    Thương mại hóa 1Stream v1.0: onboarding sản phẩm, auto-reply có
    human-in-the-loop, dashboard live-session, thống kê bình luận/lead nóng, xuất
    kết quả phiên live và cấu hình gói Free/Growth/Pro. &
    Tăng trưởng bằng early-bird pricing, case study từ khách hàng beta, demo
    trước/sau về chi phí nhân sự và số bình luận được xử lý. Theo dõi CAC, churn,
    GPU cost/job, chi phí support và tỷ lệ nâng cấp. &
    Có khách hàng trả phí đầu tiên, case study định lượng, dashboard theo dõi
    phiên live và dữ liệu thực tế để cải thiện model/prompt. \\ \hline

    \textbf{Giai đoạn 4}\newline \textbf{05--12/2027}\newline \textbf{Scale-Up \& Optimization} &
    Mở rộng ra các thành phố lớn, chuẩn hóa hạ tầng, quy trình triển khai và mô
    hình vận hành cho nhiều ngành hàng có lưu lượng bình luận cao. &
    Nâng cấp kiến trúc production: hàng đợi job bền vững, object storage,
    Redis/session store, PostgreSQL/migration, observability, autoscaling cho AI
    workers và tích hợp ERP/CRM/order service cho khách hàng tổ chức. &
    Mở rộng sang Hà Nội, Đà Nẵng và các cụm seller lớn; phát triển kênh B2B
    reseller/agency, hợp tác cộng đồng seller và thử nghiệm gói Enterprise. Gọi
    vốn seed để mở rộng hạ tầng, sales và customer success. &
    Hệ thống chịu được nhiều phiên live đồng thời, có SLA nội bộ, dữ liệu vận
    hành theo ngành hàng và kênh phân phối đủ mạnh để mở rộng ngoài TP.HCM. \\
    \hline

\end{longtable}
\end{center}

\subsubsection*{Các mốc ưu tiên gần hạn}

\begin{enumerate}
    \item \textbf{Ổn định lõi AI Sales Agent:} chuẩn hóa Knowledge Database, tăng
    độ chính xác phân loại bình luận, thêm cơ chế từ chối khi thiếu dữ liệu và
    lưu lại prompt/model version để đánh giá chất lượng.
    \item \textbf{Hoàn thiện workflow livestream:} kết nối rõ hơn giữa dashboard,
    bình luận realtime, auto-reply, hot lead, đơn hàng và báo cáo sau phiên live.
    \item \textbf{Kiểm soát chi phí AI video:} tách tác vụ nặng sang queue, batch
    rendering, lưu trữ kết quả dùng lại và chỉ dùng GPU/API model khi thật sự cần.
    \item \textbf{Chuẩn bị thương mại hóa:} thiết kế gói Standard/Growth/Pro,
    quy trình onboarding dữ liệu sản phẩm và bộ KPI để chứng minh lợi ích giảm
    chi phí nhân sự, tăng tốc độ phản hồi và tăng khả năng chốt đơn.
\end{enumerate}

